\documentclass[10pt,a4paper]{amsart}
\usepackage[margin=19mm]{geometry}
\usepackage{amsmath,amssymb,mathtools}
\usepackage[scr=rsfs,cal=euler]{mathalfa}
\usepackage{xfrac}
\usepackage[unicode,hidelinks,bookmarks=false]{hyperref}
\newcommand{\ie}{i.e.\ }
\newcommand{\cf}{cf.\ }
\newcommand{\resp}{respectively\ }
\newcommand{\Subsection}{\subsection}
\providecommand{\nobreakdash}{\nobreak\hbox{-}}
\setlength{\emergencystretch}{2em}
\multlinegap0pt
\usepackage{amssymb}
\newtheorem{definition}{Definition}
\newtheorem{proposition}{Proposition}
\newtheorem{lemma}{Lemma}
\newcommand\PP{\mathcal{P}}
\title{On groups with Schottky set boundary}
\author{Peter Ha\"issinsky, Luisa Paoluzzi,, Genevieve Walsh}
\date{Source: arXiv:2309.08735v3 (CC BY 4.0)}
\begin{document}
\maketitle

\noindent\emph{Source and licence.} On groups with Schottky set boundary. Version arXiv:2309.08735v3 (CC BY 4.0), distributed by arXiv under the Creative Commons Attribution 4.0 International licence, http://creativecommons.org/licenses/by/4.0/, as recorded in the arXiv licence metadata for this exact version and shown on the abstract page https://arxiv.org/abs/2309.08735v3. Full licence notice: \texttt{licenses/arxiv-2309.08735-CC-BY-4.0.txt}.\par\smallskip

\noindent\textit{Editorial scope.} Counted unit: the article's lemma lma:parabolicvsfuchs (source0.tex lines 535--537). The article's complete original proof of it is delivered with it (source0.tex lines 539--570, one \texttt{proof} environment of 32 lines). No mathematics is written, repaired or re-derived: the statement and the proof are the article's own lines, in the article's own notation, and no retained line is substituted or re-flowed.  Retained uncounted as the source-local context the proof needs: lines 226--230; lines 336--339.  Nothing was omitted from any retained span, so the package carries no omission record.  No retained line contains a citation, so the wrapper carries no bibliography and no bibliography entry.  No retained line contains a figure environment, an image-inclusion command or drawing code, so no image is delivered and the package declares no asset dependency.  Editorial formatting only, in addition: an AMS \texttt{amsart} preamble that restates the article's own declarations and macros used by the retained lines, namely the environments \texttt{definition}, \texttt{proposition}, \texttt{lemma} on separate counters, together with the standard packages those lines need.  The retained environments print in this document as Definition 1, Proposition 1, Lemma 1 in order of appearance, whereas the article prints the same items with the numbering of its own full document.  The complete native source of the article as distributed in the arXiv e-print for this exact version is bundled unchanged as \texttt{sources/147105/source0.tex}.
\begin{definition}[Null sequences and $E$-sets]\label{def:nses}
Given a compact metric space $Z$, a \emph{null-sequence} is a collection of subsets $\mathcal{C}$ such that, for any $\delta >0$, the collection $\mathcal{C}$ contains at most finitely many elements of diameter at least $\delta$.

An \emph{$E$-set} is a connected compact subset of the sphere $S^2$ such that the collection of connected components of its complement is a null-sequence.
\end{definition}

\begin{proposition}\label{prop:infends} The set of components of the Bowditch boundary of a relatively hyperbolic group $(G,\PP)$ 
forms a null-sequence. Moreover, for each 
component containing at least two points, the stabilizer is hyperbolic relative to conjugates of the original
peripheral subgroups $\PP$. \end{proposition}

\begin{lemma}\label{lma:parabolicvsfuchs}  Let $(G, \PP)$ be a geometrically finite, non elementary, convergence group on $S^2$ with limit set $\Lambda$.  Let $\Omega$ be a simply connected component of $S^2 \setminus \Lambda$ and {$f:\mathbb{D}\to \Omega$ the homeomorphism that conjugates the stabilizer $H$ of $\Omega$} to a Fuchsian group, and $h:\overline{\mathbb{D}}\to\overline{\Omega}$ the extension of this homeomorphism as above. Let $p\in\partial\Omega$ be a parabolic point with stabilizer $P$, and set {$Q= f^{-1}\circ(P\cap H)\circ f$}. 
Then the limit set $\Lambda_Q$ is exactly the non-empty set  $h^{-1}(\{p\})$.
\end{lemma}

\begin{proof}  First use Proposition \ref{prop:infends} to reduce to the case when the boundary is connected. 
Let $K\subset S^2\setminus\{p\}$ be a compact subset containing a fundamental domain for the action
of $P$ on $\Lambda\setminus\{p\}$.  We may find
$g\in P$ such that $g(\partial\Omega)\cap K\ne \emptyset$. 
Thus, $g(\Omega)$ contains in its closure
the point $p$ and at least one point of $K$. Such components form a finite set since $\Lambda$ is an $E$-set (Def. \ref{def:nses}). 
By considering a sequence of points in $\partial\Omega$ tending to $p$, we may pick an infinite sequence $(g_n)$ in $P$ that
maps $\Omega$ to components whose closures intersect both $\{p\}$ and $K$.
As there are only finitely many of them,  we may assume that $g_n(\Omega)=V$ for a fixed component 
$V$,
and all $n\ge 1$. Therefore, $(g_1^{-1}g_n)_n$
is an infinite collection of elements of $H\cap P$ which proves that $\Lambda_Q$ is not empty.

Since $hQ\subset Ph$, it follows that $h^{-1}(\{p\})$ is $Q$-invariant and compact, hence it contains
$\Lambda_Q$ which by definition is the minimal  compact invariant subset under the action of $Q$. 

The equality will follow from the fact that $p$ is a bounded parabolic point. We first rule out the case that $h^{-1}(\{p\})$ contains
an interval. If this were the case, then it would be the whole circle by invariance so that we would have $H=P$; this
contradicts that $\Lambda_P=\{p\}$ and $\Lambda_H=\partial\Omega$. 
Therefore, $h^{-1}(\{p\})$ is nowhere dense in $\mathbb{S}^1$.

Let $\Omega_1,\ldots \Omega_k$, be the $P$-translates
of $\Omega$ whose closures intersect both $\{p\}$ and $K$ and let us fix $g_1,\ldots g_k\in P$ such that
$g_j(\Omega_j)= \Omega$. Set $L= h^{-1}(\cup_{1\le j\le k} g_j(K))$.  This is a compact subset of $\overline{\mathbb{D}}$ disjoint from $h^{-1}(\{p\})$, hence from $\Lambda_Q$.
Let $x\in h^{-1}(p)$.  
We want to prove that the action of $Q$ is not equicontinuous at $x$.
With that in mind,
pick a point $y\in\mathbb{S}^1\setminus h^{-1}(\{p\})$ arbitrarily close to $x$. 
Note that $h(y)\in\partial\Omega\setminus\{p\}$ and since $K$ is a fundamental domain, we may find $g\in P$ and $j\in\{1,\ldots, k\}$ so that $g(h(y))\in K\cap\partial\Omega_j$.
It follows that $g_jg\in (H\cap P)$ so that we may find $q= h^{-1}g_jgh\in Q$ with $q(y)\in L$. This implies that $x\in\Lambda_Q$. 
Indeed, considering now a sequence $(y_n)$ in $\mathbb{S}^1\setminus h^{-1}(\{p\})$ tending to $x$, we obtain in this way a sequence $(q_n)$ in $Q$ such that $(q_n(x), q_n(y_n))\in h^{-1}(\{p\})\times L$: as $L$ and  $h^{-1}(\{p\})$ are disjoint compact subsets, $(q_n)_n$ cannot be equicontinuous at $x$.
\end{proof}

\end{document}
